\section{Recap on Follow-The-Regularised-Leader}
This section will consist of a brief recap of the Follow-the-Regularised-Leader (FoReL) back in Week 7. FoRel is a simple modification of the Follow-The-Leader (FTL) learning rule.  In FTL, in each iteration, we select the decision vector (or the next iterate) which minimises the loss on all past iterations:
\begin{define}[Follow-The-Leader]
    The Follow-The-Leader (FTL) learning rule is as such: we select the next iterate my minimising the loss on all past rounds. That is for a convex set $S$,
    \begin{equation}
        \forall t\; x_t = \arg \min_{x \in S} \left\{\sum_{i = 1}^{t-1} f_i(x)\right\}
    \end{equation}
\end{define}
In FoRel, we add a regularisation term $h$. The goal of the regularisation term is to stabilise the solution.
\begin{define}[Follow-The-Regularised-Leader]\label{defn:forel} The Follow-The-Regularised-Leader (FoREL) learning rule extends FTL by adding a regularisation term $h: S \to \mathbb R$.
    \begin{equation}
        \forall t\; x_t = \arg \min_{x \in S} \left\{\sum_{i = 1}^{t-1} f_i(x) + h(x)\right\}
    \end{equation}
\end{define}

\subsection{Mathematical Preliminaries}
We also recap some of the important definitions that will be important in helping us understand the tools used for the analysis of regret bounds.

\begin{define}[L-Lipschitz]
    A function $f$ is called $L$-Lipschitz over a set $S$ with respect to a norm $||\cdot||$ if for all $x, x' \in S$, we have that \begin{equation*}
        |f(x) - f(x')| \leq L ||x - x'||
    \end{equation*}
\end{define}

\begin{define}[$\sigma$-strongly convex]\label{defn:sigma_strong_convexity}
    A function $f$ that is $\sigma$-strongly convex over a set $S$ with respect to norm $||\cdot ||$ if for any $x' \in S$, we have that \begin{equation*}
        \forall z \in \partial f(x'),\;\forall x \in S, \; f(x) \geq f(x') + \langle z, x - x' \rangle + \frac{\sigma}{2} || x - x' ||^2
    \end{equation*}
    where $\partial f(x')$ is the set of sub-gradients of $f$ at $x'$. Furthermore, if $x^*$ minimises $f$ over $S$, then for all $x' \in S$
    \begin{equation*}
        f(x') - f(x^*) \ge \frac{\sigma}{2}||x' - x^*||^2
    \end{equation*}
\end{define}

We also define something that might be unfamiliar to the reader which is called the \textit{dual norm}, denoted by $||\cdot||_*$
\begin{define}
    Let $||\cdot||$ be a norm. The dual norm of $||\cdot||$ is a function $||\cdot||_*$ defined as
    \begin{equation}
        ||y||_* := \sup\{\langle x, y \rangle: ||x|| \leq 1\}
    \end{equation}
\end{define}
The dual norm of the $\ell_2$-norm is again the $\ell_2$-norm (the $\ell_2$-norm is therefore \textit{self-dual}). In general, the dual for the $\ell_p$-norm is the $\ell_q$-norm, where $\frac{1}{p} + \frac{1}{q} = 1$. We also give a lemma for a convex function which is $L$-Lipschitz and how it relates to the dual norm.
\begin{lemma} Let $f: S \to \mathbb R$ be a convex function. $f$ is $L$-Lipschitz over $S$ with respect to a norm $||\cdot||$ if and only if for all $x'\in S$ and $z \in \partial f(x')$, we have $||z||_* \leq L$.
\end{lemma}
\subsection{Regret Bound Results for FoReL}
We give a proof sketch on deriving the regret bound for FoReL. From Week 7's notes, the regret for FoReL can be written as the following when evaluated against a fixed value $x$ as such:
\begin{equation}\label{eqn:regret_forel}
    \text{Regret}_T(x) := \sum_{t=1}^T f_t(x_t) - \sum_{t=1}^T f_t(x)
\end{equation}
To prove the bound for the above expression, recall also that we have the following lemma:
\begin{lemma}\label{lemma:iterates_forel}
    Let $x_1, x_2, \ldots$ be the iterates due to FoReL using the $\sigma$-strongly convex regulariser $h$. Then for all $x \in S$, we have that \begin{equation}\label{eqn:lemma_forel}
        \sum_{t=1}^T (f_t(x_t) - f_t(x)) \leq h(x) - h(x_1) + \sum_{t=1}^T (f_t(x_t) - f_t(x_{t+1}))
    \end{equation}
\end{lemma}
Furthermore, it is also shown that if for all $t$, if $f_t$ is $L_t$-Lipschitz with respect to $||\cdot||$, then we further have that 
\begin{equation}\label{eqn:lemma_forel_lipschitz}
    f_t(x_t) - f_t(x_{t+1}) \leq L_t ||x_t - x_{t+1}|| \leq \frac{L_t^2}{\sigma}
\end{equation} 
Again, interested readers may refer to Week 7's notes or \cite{shalev2025online} for the full derivation. Combining Equations~\ref{eqn:lemma_forel} and \ref{eqn:lemma_forel_lipschitz} we have that the regret bound for FoReL using the regulariser $h$ is as follows:
\begin{equation}\label{eqn:forel_regre_with_min}
    \text{Regret}_T(x) \leq h(x) - h(x_1) + \frac{TL^2}{\sigma} \leq h(x) - \min_{v \in S} h(v) + \frac{TL^2}{\sigma}
\end{equation} where $L$ is such that $\frac{1}{T}\sum_{t=1}^T L_t^2 \leq L^2$.

\subsection{The Problem with FoReL}
The main issue (or a disadvantage) is that for FoReL, we need solve an optimisation problem \textit{at every iteration}, which is computationally inefficient and might be extremely complex. This leads us to a very natural question: \textit{Can we achieve a similar regret bound with a simpler update rule?} It turns out that we indeed can achieve a much simpler update rule without worsening the regret bound by using the Online Mirror Descent (OMD) family of algorithms.